\documentclass[10pt, aspectratio=169]{beamer}

\usetheme{default}
\usecolortheme{dove}
\setbeamertemplate{navigation symbols}{}
\setbeamerfont{title}{series=\bfseries}
\setbeamerfont{frametitle}{series=\bfseries}
\setbeamercolor{frametitle}{fg=black}

\usepackage{booktabs}
\usepackage{graphicx}
\usepackage{amsmath}
\usepackage[T1]{fontenc}
\usepackage{lmodern}
\usepackage{hyperref}
\hypersetup{colorlinks=true, urlcolor=blue!50!black, linkcolor=black}

\newcommand{\paperlink}[2]{\href{#1}{\textcolor{blue!50!black}{\footnotesize [#2]}}}

\AtBeginSection[]{%
  \begin{frame}[plain]
    \vfill\centering{\usebeamerfont{title}\insertsectionhead}\vfill
  \end{frame}
}

\title{Global impact of rainfall on mortality: drought channel}
\author{Marine and Tom}
\date{\today}

\begin{document}

\begin{frame}[plain]\titlepage\end{frame}


\begin{frame}{Purpose}
  \textbf{Goal:} estimate the impact of rainfall on mortality, globally.
  \vspace{1em}
  \begin{itemize}
    \item \textbf{Challenge:} the impacts of rainfall on mortality are highly heterogeneous.
    \item Droughts kill --- but so does extreme rainfall. 
    \item Different functions of rainfall impact mortality through different mechanisms, at different temporal scales, in different places.
  \end{itemize}
  \vspace{1em}
  A single pooled response function may average effects toward zero.
\end{frame}

\begin{frame}{Plan}
  \vfill
  \begin{columns}[c]
    \begin{column}{0.12\textwidth}
      \centering{\fontsize{10}{10}\selectfont\color{black!50}\bfseries Step 1}
    \end{column}
    \begin{column}{0.8\textwidth}
      {\large\textbf{Paradigms}}\\[0.4em]
      Build a series of mechanism-specific paradigms, informed by the literature, showing empirically how different mechanisms operate in different parts of the sample.
    \end{column}
  \end{columns}
  \vspace{1.2em}
  {\color{black!20}\hrule}
  \vspace{1.2em}
  \begin{columns}[c]
    \begin{column}{0.12\textwidth}
      \centering{\fontsize{10}{10}\selectfont\color{black!50}\bfseries Step 2}
    \end{column}
    \begin{column}{0.8\textwidth}
      {\large\textbf{Model selection}}\\[0.4em]
      Use machine learning to allow for rich heterogeneity everywhere, informed by the paradigms from Stage 1.
    \end{column}
  \end{columns}
  \vfill
\end{frame}

\begin{frame}{Current status: paradigm 1 --- extreme rainfall in cities}
  \begin{itemize}
    \item Bearpark et al.\ (2025, \emph{Nature}): extreme daily rainfall causes large mortality impacts in Mumbai.
    \item Suggestive evidence of the same pattern in our sample, when restricting to admin units that resemble Mumbai (dense, urban, flood-exposed).
  \end{itemize}
  \vspace{2em}
  Paradigm 1: Cities may carry large negative effects of \emph{extreme daily} rainfall.
\end{frame}

\begin{frame}{Paradigm 2: drought}
  \vfill
  \centering
  {\Large How should we measure drought,\\[0.3em]
  and where should it matter for mortality?}
  \vfill
\end{frame}

%============================================================================%
\section{Literature}
%============================================================================%

\begin{frame}{The Economics of Drought}
    
  \textbf{Core insight: the rainfall deficit is only the beginning of the story.}
  Whether it becomes an actual shortage depends on storage, groundwater,
  irrigation, conveyance, prices, allocation rules, markets, \ldots
  \vspace{0.6em}
  \begin{itemize}
    \item \textbf{drought} $\neq$ \textbf{water scarcity} $\neq$ \textbf{water shortage}.
    \begin{itemize}
      \item A drought can raise scarcity without creating a shortage that households or farms experience.
      \item Conversely: high demand or poor allocation can create a shortage with no exceptional meteorological drought.
    \end{itemize}
  \end{itemize}
  \vspace{1em}
  A useful frame for choosing heterogeneity dimensions and mechanisms.

  \vfill
  \paperlink{https://doi.org/10.1146/annurev-resource-101623-100253}{Edwards, Sanchez \& Sekhri 2024}\\
\end{frame}

\begin{frame}{Water Scarcity and Birth Outcomes in the Brazilian Semiarid}
  Rainfall deficits during pregnancy worsen birth outcomes and raise infant mortality in poor semiarid places, mainly through intestinal infections and malnutrition.
  \vspace{0.6em}
  \begin{itemize}
    \item \textbf{Key heterogeneity:} mortality effects are larger where piped-water and sanitation coverage are low (and for youngest age group). 
    \item Measures rainfall during 
    \item Null impacts for agriculture
  \end{itemize}

  \vfill
    \paperlink{https://doi.org/10.1016/j.jdeveco.2014.10.003}{Rocha \& Soares 2015}
\end{frame}

\begin{frame}{Global WASH mapping }
  \textbf{A data source for heterogeneity:} annual 5\,km estimates of water-source
  and sanitation categories, 90 LMICs, 2000--2017.
  \vspace{0.6em}
  \begin{itemize}
    \item Rasters plus adm1/adm2 summaries; the rasters let us aggregate to
          \emph{our own} admin boundaries.
    \item Four categories each for water source and sanitation.
          Initially keep them \textbf{separate} rather than collapsing into a
          single WASH index --- they may capture different dimensions of
          vulnerability.
  \end{itemize}
  \vspace{0.6em}
  \textbf{Open question:} let these measures vary through time, or take a fixed
  baseline / sample average?

  \vfill 
    \paperlink{https://doi.org/10.1016/S2214-109X(20)30278-3}{Local Burden of Disease WASH Collaborators 2020}
\end{frame}

\begin{frame}{Drought and mortality in urban Brazil}
  Weekly mortality in 13 large Brazilian urban areas, linked to SPEI-1.
  \vspace{0.6em}
  \begin{itemize}
    \item Pooled effects small and imprecise --- but stronger signals for
          respiratory and circulatory mortality, and among women, children and
          older adults.
    \item \textbf{Takeaway:} do not exclude cities from the drought analysis.
          Short dry periods may affect particular urban populations even when
          the average city-level relationship is weak.
    \item Mechanisms not identified; candidates discussed: water-system
          stress, dust / air pollution, dehydration, heat, socioeconomic
          vulnerability. A good starting list.
    \item \textbf{Measurement issue:} they use city-center exposure, but drought
          may hit a city's \emph{source watershed or reservoir system}. How to
          measure that? (Proxy from James's hydrology work?)
  \end{itemize}
  \vfill
    \paperlink{https://doi.org/10.1029/2021GH000534}{Salvador et al.\ 2022}
\end{frame}

%============================================================================%
\section{Defining and Measuring Drought}
%============================================================================%

\begin{frame}{What is drought?}
  \textbf{Working definition:} a rainfall shortage that is \emph{unusual} for a
  place and time of year, and that \emph{persists} long enough to matter.
  \vspace{0.6em}

  Three features from the literature
  \paperlink{https://doi.org/10.1080/02508068508686328}{Wilhite \& Glantz 1985}
  \paperlink{https://doi.org/10.1029/WR016i002p00297}{Dracup, Lee \& Paulson 1980}
  \paperlink{https://www.droughtmanagement.info/find/guidelines-tools/handbook-drought-indicators-and-indices/}{WMO/GWP 2016}:
  \begin{itemize}
    \item \textbf{Relative:} 20\,mm can be extremely dry in one place/month and
          completely normal in another. A normal dry season is not drought.
    \item \textbf{Persistent:} one dry day, week or even month is a rainfall
          shock; drought is a deficit that continues or accumulates.
    \item \textbf{Propagates:} precipitation deficit $\rightarrow$ low soil
          moisture $\rightarrow$ low streamflow / storage $\rightarrow$ crop
          stress $\rightarrow$ experienced water or food shortage. Related, but
          \emph{not} the same exposure, nor the same time scale.
  \end{itemize}
  \vspace{0.4em}
  Four categories \paperlink{https://doi.org/10.1080/02508068508686328}{Wilhite \& Glantz 1985}:
  \textbf{meteorological} $\cdot$ agricultural $\cdot$ hydrological $\cdot$
  socioeconomic. Our common cross-country exposure is \textbf{meteorological};
  the others are mechanism-specific measures or consequences (agricultural
  drought may need a growing-season definition --- to discuss).
\end{frame}

\begin{frame}{Proposed construction: from admin-month rainfall to episodes}
  Dimensions worth preserving
  \paperlink{https://doi.org/10.1029/WR016i002p00297}{Dracup, Lee \& Paulson 1980}:
  onset $\cdot$ duration $\cdot$ intensity $\cdot$ cumulative severity $\cdot$
  worst month (``nadir'') $\cdot$ recovery. \emph{Which do we need?}
  \vspace{0.6em}
  \begin{enumerate}
    \item \textbf{Admin-month rainfall.} Sum daily admin-level rainfall to
          monthly totals (already processed; pop-weighted as default?).
    \item \textbf{Fixed reference period.} Same across the panel --- but which?
          Projection-system ``historical climate'' (1981--2001? 1981--2010?).
          Country-specific baselines would change what the percentile means.
          Alternative reference periods as sensitivity checks.
          \emph{Tom: rolling average} --- needs precip data beyond mortality
          coverage years.
    \item \textbf{Local rainfall percentile.} Compare each admin-month with the
          historical distribution for the \emph{same admin, same calendar
          month} (January vs.\ Januaries). Solves cross-climate comparability;
          retain totals and mm deficit alongside.
    \item \textbf{Join months into episodes.} First rule: $\leq$ 20th local
          percentile for $\geq$ 2 consecutive months (test alternatives).
          Per episode keep onset, duration, worst percentile, cumulative
          percentile and mm deficits.
  \end{enumerate}
\end{frame}

\begin{frame}{Standard indices, and getting to the annual panel}
  \textbf{Standard indices as complements?}
  \begin{itemize}
    \item \textbf{SPI-3 / 6 / 12} from monthly precip
          \paperlink{https://climate.colostate.edu/pdfs/relationshipofdroughtfrequency.pdf}{McKee et al.\ 1993}
          \paperlink{https://library.wmo.int/idurl/4/39629}{WMO 2012}:
          widely used, comparable across climates; less intuitive than direct
          duration/severity. Windows map to mechanisms: SPI-3 seasonal /
          agriculture, SPI-12 longer-term water supply.
          Constructed at adm-unit level (settled).
    \item \textbf{SPEI} as robustness
          \paperlink{https://doi.org/10.1175/2009JCLI2909.1}{Vicente-Serrano et al.\ 2010}:
          adds evapotranspiration --- but then partly combines low precip with
          temperature.
  \end{itemize}
  \vspace{0.6em}
  \textbf{Monthly episodes $\rightarrow$ annual mortality panel.} Identify
  episodes in monthly data first, then translate to admin-year measures:
  \begin{itemize}
    \item months in qualifying drought during the year; maximum episode duration
    \item cumulative percentile and/or mm deficit; worst monthly percentile
    \item minimum SPI-3 / SPI-6 / SPI-12
  \end{itemize}
\end{frame}

%============================================================================%
\section{Next Steps}
%============================================================================%

\begin{frame}{First mechanism: household water and sanitation}
  The clearest place to start --- we have a plausible cross-country data source.
  \vspace{0.6em}

  \textbf{The story:} persistent low rainfall reduces water quantity or
  reliability $\rightarrow$ households store water or switch sources
  $\rightarrow$ hygiene and sanitation worsen $\rightarrow$ infection,
  dehydration, nutrition risks rise. Expect the strongest effects where
  piped-water and sanitation access are low, surface/unimproved water is
  common, and among infants and young children.
  \vspace{0.6em}

  \textbf{What we can realistically do:}
  \begin{itemize}
    \item Aggregate the LBD WASH estimates to our admin units.
    \item Build baseline (or time-varying) shares: piped water,
          surface/unimproved water, sewer/septic, \ldots
    \item Interact these with drought --- separately at first, not as one index.
    \item Focus on the youngest age group and enteric/malnutrition causes where
          available (only a few countries in our sample).
  \end{itemize}
\end{frame}

\end{document}
